\documentclass[11pt,a4paper]{article}
\usepackage[T1]{fontenc}
\usepackage[utf8]{inputenc}
\usepackage{lmodern,microtype}
\usepackage[margin=26mm,headheight=15pt]{geometry}
\usepackage{amsmath,amssymb,amsthm,mathtools}
\usepackage{booktabs,array,longtable,enumitem}
\usepackage[dvipsnames]{xcolor}
\usepackage{fancyhdr,xurl}
\usepackage[colorlinks=true,linkcolor=MidnightBlue,citecolor=MidnightBlue,urlcolor=MidnightBlue,unicode]{hyperref}
\usepackage{bookmark,listings}
\hypersetup{pdftitle={Sharp negative-direction summability for countable exponential feedback},pdfauthor={Research article prepared with ChatGPT for Vladimir Reshetnikov}}
\pagestyle{fancy}\fancyhf{}
\fancyhead[L]{\small Countable exponential feedback}
\fancyhead[R]{\small Sectorial asymptotics and summability}
\fancyfoot[C]{\thepage}\renewcommand{\headrulewidth}{0.3pt}
\setlength{\parskip}{3pt}\setlength{\emergencystretch}{2em}
\setlist{itemsep=4pt,topsep=6pt,leftmargin=1.7em}
\numberwithin{equation}{section}
\newtheorem{theorem}{Theorem}[section]
\newtheorem{proposition}[theorem]{Proposition}
\newtheorem{lemma}[theorem]{Lemma}
\newtheorem{corollary}[theorem]{Corollary}
\theoremstyle{definition}\newtheorem{definition}[theorem]{Definition}
\theoremstyle{remark}\newtheorem{remark}[theorem]{Remark}
\newcommand{\N}{\mathbb N}\newcommand{\C}{\mathbb C}\newcommand{\R}{\mathbb R}
\newcommand{\Uhat}{\widehat U}\newcommand{\Qhat}{\widehat Q}\newcommand{\Vhat}{\widehat V}
\newcommand{\M}{\mathbb M}\newcommand{\CM}{\mathcal C_{\M}}
\newcommand{\dd}{\,\mathrm d}\newcommand{\ee}{\mathrm e}
\DeclareMathOperator{\Rea}{Re}
\lstset{basicstyle=\ttfamily\small,breaklines=true,columns=fullflexible,frame=single,showstringspaces=false}
% ed. (2026-09-29): editorial-note environment for ProveIt cross-references;
% it uses no counter, so no author numbering changes.
\newenvironment{ednote}{\par\smallskip\begin{quote}\small\raggedright\textbf{Editorial note (ProveIt, 2026-09-29).}\ }{\end{quote}\smallskip}
\begin{document}
% ed. (2026-09-29): no page anchor on the title page, removing the duplicate
% hyperref destination page.1 (the contents page is also numbered 1).
\hypersetup{pageanchor=false}
\begin{titlepage}\centering
\vspace*{12mm}
{\small\scshape Research article\par}\vspace{9mm}
{\LARGE\bfseries Sharp negative-direction summability\par
for countable exponential feedback\par}\vspace{6mm}
{\Large A coefficient-to-remainder principle,\par
optimal logarithmic scales, and sectorial inversion\par}\vspace{10mm}
{\large Prepared for Vladimir Reshetnikov\par}\vspace{3mm}
29 September 2026\vspace{10mm}
\begin{minipage}{0.93\textwidth}
\begin{abstract}
Consider the formal and sectorial equation
\[
 U(q)=\sum_{j\ge1}q^j\exp(\lambda_jU(q)),\qquad \lambda_j\ge0.
\]
A recent ProveIt draft establishes coefficient regularity and a left-sector
realization but explicitly leaves uniform Gevrey remainders unproved.
We close that gap for a broad class of weights. The change of unknown
$U=qV$ gives a fixed dependent-variable disk about $V=1$. An exact
exponential-remainder estimate links the kernel to a positive subfamily
of the Lagrange coefficients. We prove that the canonical sectorial
solution has strong $\mathbb M$-asymptotics if and only if its formal
coefficients belong to $\mathcal C_{\mathbb M}$, under explicit elementary
weight hypotheses. In particular, strong Gevrey-$s$ asymptotics are
equivalent to $\lambda_j=O((j\log(j+1))^{s+1})$ for every $s>0$.
The classical sectorial characterization of summability then gives a
sharp application: for every $k>1$, negative-direction $k$-summability
holds if and only if $\lambda_j=O((j\log(j+1))^{1+1/k})$.
For $\lambda_j\asymp j^p(\log(\ee+j))^\beta$, $p>1$, we determine
the optimal factorial--logarithmic weight, up to geometric factors, and
prove actual truncation errors on the scale
$\exp[-c|q|^{-1/(p-1)}(\log(1/|q|))^{(p-\beta)/(p-1)}]$.
The results extend to the sectorial compositional inverse. Detailed
proofs, error certificates, exact checks, limitations, and further
research questions are included.
\end{abstract}
\end{minipage}\vfill
\begin{minipage}{0.93\textwidth}\small
\textbf{Status.} These are conventional mathematical proofs, not
Lean-verified theorems. The model-specific coefficient-to-remainder and
summability conclusions are the proposed contribution. General weighted
implicit-function methods and the classical summability theorem are not
claimed as inventions. Global novelty and priority have not been
certified. No named longstanding conjecture is claimed solved.
\end{minipage}
\end{titlepage}
% ed. (2026-09-29): page anchors on again after the title page.
\hypersetup{pageanchor=true}
\tableofcontents\clearpage

\section{The research gap and the main conclusions}
\label{sec:intro}
\subsection{A precise continuation of the repository}
The ProveIt snapshot used here is commit
\begin{center}\small\texttt{0c973d8f5e7ef4550f4df1bf486d890037fd9f14}.\end{center}
Its transseries documentation is in
\path{Analysis/Transseries/docs/series-and-transseries/}.
The directory README distinguishes the consolidated volumes from nineteen
unmerged research arrivals dated 29 September 2026, and warns that those
arrivals have not been reviewed claim by claim or formalized \cite{repo}.

The immediate predecessor is \emph{A sharp regularity classification for
countable exponential-feedback transseries}, in
\path{Exponential_Feedback_Regularity_Classification/} \cite{feedback}.
It studies
\begin{equation}
 \Uhat(q)=\sum_{j\ge1}q^j\exp(\lambda_j\Uhat(q)),\qquad
 \Uhat(q)=\sum_{n\ge1}u_nq^n,\quad \lambda_j\ge0.
 \label{eq:model}
\end{equation}
Its README and sectorial-realization proof explicitly distinguish
Poincar\'e asymptotics from uniform Gevrey remainder estimates. For each
fixed $N$, the assertion $R_N(q)=O_N(q^N)$ alone says nothing useful about
how the constant grows with $N$.

The neighboring \emph{Finite-Core Universality and Sharp Large Order}
draft develops multiplicative coefficient asymptotics and finite-core
condensation \cite{finitecore}. We do not repeat its saddle analysis.
The present question is to control the actual error of a sectorial
solution and identify which summation procedures reconstruct it.
% ed. (2026-09-29): missing citation: the negative-ray package was in this article's snapshot but is not cited.
\begin{ednote}
The snapshot used here already contained a third package, \nolinkurl{Negative_Ray_Summation_Exponential_Feedback/}~\cite{ed:negativeray}, which the directory README described
and which this article does not cite. It proves, for $s=1$ and unit
amplitudes, that $\Uhat$ and its compositional inverse are \emph{finely}
(fixed-width half-strip) Borel--Laplace summable in direction $\pi$
exactly when $\lambda_j=O((j\log j)^2)$ (its Theorem~\texttt{thm:threshold}),
including the quadratic model, by inverting the kernel before taking Borel
transforms; it also continues the inverse Borel transform explicitly and,
for $\lambda_j=j^2$, bounds the inverse remainder uniformly on a closed
left half-disc (its Theorems~\texttt{thm:borel} and \texttt{thm:remainder}). The two
articles are complementary: this one covers every $s>0$ and sectorial
$k$-summability for $k>1$; that one covers the fine case $k=1$; both prove
the Gevrey-$1$ class equivalence. The ordinary (angular) case $k=1$ for
$\lambda_j=j^2$ is settled negatively by the later package \nolinkurl{Natural_Boundaries_Quadratic_Exponential_Feedback/}~\cite{ed:naturalboundaries}.
All three packages are in
\nolinkurl{Analysis/Transseries/docs/series-and-transseries/}.
\end{ednote}

\subsection{Principal statements}
\begin{theorem}[Sharp sectorial Gevrey and summability classification]
\label{thm:headline}
Let $(\lambda_j)$ be any sequence of finite nonnegative real numbers.
There is a canonical holomorphic germ $U$ on a left sectorial region of
opening $\pi$, characterized by \eqref{eq:model} and $U=q+O(q^2)$.
For every $s>0$, the following conditions are equivalent:
\begin{enumerate}[label=(\roman*)]
\item $\lambda_j=O((j\log(j+1))^{s+1})$;
\item $|u_n|\le CA^n(n!)^s$ for some $C,A>0$;
\item on every proper left sector, after choosing a sufficiently small
radius, there are $C,A>0$ such that for all $N\ge1$,
\begin{equation}
 \left|U(q)-\sum_{n=1}^{N-1}u_nq^n\right|
 \le CA^N(N!)^s|q|^N.
 \label{eq:headline-remainder}
\end{equation}
\end{enumerate}
For every $k>1$, these conditions with $s=1/k$ are also equivalent to
$\Uhat$ being $k$-summable in direction $\pi$. Its sum is $U$.
The same equivalences hold for the formal and sectorial compositional
inverses of $U$.
\end{theorem}
The summability conclusion deliberately assumes $k>1$; the case $k=1$
is not settled here. The remainder assertion holds for every $s>0$.
% ed. (2026-09-29): scope of "Sharp" in the title: the summability iff holds for k > 1 only.
\begin{ednote}
Scope: the word ``Sharp'' in the title refers to these equivalences. For
summability they hold for $k>1$ only, that is for subquadratic slope
classes; the quadratic model, which is the regularity article's directional
question, is not decided here. Its fine $k=1$ summability is proved in
\nolinkurl{Negative_Ray_Summation_Exponential_Feedback/}~\cite{ed:negativeray}, and its angular $1$-summability fails for $\lambda_j=j^2$
by \nolinkurl{Natural_Boundaries_Quadratic_Exponential_Feedback/}~\cite{ed:naturalboundaries}.
\end{ednote}

\begin{theorem}[Optimal logarithmic scale]
\label{thm:logheadline}
Suppose, for some $p>1$, $\beta\in\R$, and $c,C>0$, that
\begin{equation}
 cj^p(\log(\ee+j))^\beta\le\lambda_j
 \le Cj^p(\log(\ee+j))^\beta
 \quad\text{for all sufficiently large }j.
 \label{eq:twosidedslopes}
\end{equation}
Set $s=p-1$, $\tau=\beta-p$, and, up to a finite initial quotient
adjustment, define
\begin{equation}
 M_n=(n!)^s\prod_{\ell=1}^n(\log(\ee+\ell))^\tau.
 \label{eq:logweight}
\end{equation}
Then $U$ and its sectorial inverse have strong $\M$-asymptotic expansions.
This weight cannot be replaced by $M_{s,\tau-\varepsilon}$,
$\varepsilon>0$, or by $M_{s',\tau'}$ with $0<s'<s$ and fixed $\tau'$.
On every proper left sector there is a truncation order $N=N(|q|)$ giving
\begin{equation}
 \left|U(q)-\sum_{n<N}u_nq^n\right|
 \le C_1\exp\!\left[-c_1|q|^{-1/s}
                        (\log(1/|q|))^{-\tau/s}\right].
 \label{eq:optimized-headline}
\end{equation}
The inverse has the same error scale, with possibly different constants.
If $1<p<2$, both series are also $\M$-summable in direction $\pi$ in the
standard generalized sense.
\end{theorem}
Here optimality is within the factorial--logarithmic weight scale, up to
geometric factors. It is not a lower bound on each pointwise remainder,
and does not specify the sharp numerical constant $c_1$.

\subsection{Why a change of unknown is necessary}
If $\lambda_j/j$ is unbounded, the original kernel is not analytic in a
product neighborhood of $(q,U)=(0,0)$. At fixed positive $U$, its terms
fail to tend to zero along a subsequence for every positive $q$.
An implicit-function theorem applied directly in these coordinates would
therefore lack a required hypothesis.
Instead put
\begin{equation}
 U=qV,\qquad F(q,v)=\sum_{j\ge1}q^{j-1}e^{\lambda_jqv},
 \qquad V=F(q,V).
 \label{eq:desing}
\end{equation}
On a proper left sector, $qv$ remains in the left half-plane throughout
one fixed disk about $v=1$. All exponentials are bounded by one,
independently of the slope sizes. This fixed dependent-variable disk is
the essential missing hypothesis, not a new formal inversion formula.

Gevrey asymptotic implicit-function theory is classical; for example
Nikolaev establishes a Gevrey implicit-function and resummation theorem
under appropriate uniform assumptions \cite{nikolaev}. We verify a
weighted version directly for this countable singular kernel and use
positivity to obtain an equivalence rather than only sufficiency.

\section{Conventions and admissible weights}
\label{sec:weights}
For $0<\eta<1$ and $r>0$, write
\begin{equation}
 S_{\eta,r}=\{q\in\C:0<|q|<r,\ \Rea q<-\eta|q|\}.
 \label{eq:sector}
\end{equation}
Its opening is $2\arccos\eta<\pi$. A left sectorial germ means compatible
functions on these sectors, with radii allowed to depend on $\eta$.
No radius uniform as $\eta\downarrow0$ is asserted.

\begin{definition}[Admissible weights]
A sequence $\M=(M_n)_{n\ge0}$ is admissible here if $M_0=1$, its quotients
$\mu_n=M_n/M_{n-1}$ satisfy
\begin{equation}
 1\le\mu_1\le\mu_2\le\cdots,\qquad \mu_n\longrightarrow\infty,
 \label{eq:quotients}
\end{equation}
and, for some $D,E\ge1$,
\begin{equation}
 \mu_n^n\le D^nM_n\quad(n\ge1),\qquad
 M_{n+1}\le E^{n+1}M_n\quad(n\ge0).
 \label{eq:weightconditions}
\end{equation}
\end{definition}
These are the precise hypotheses used in our proof, not a proposed
replacement for the standard definition of strong regularity. Familiar
factorial--logarithmic weights satisfy them; standard strongly regular
weights provide many examples \cite{lms}.

\begin{definition}[Coefficient and strong asymptotic classes]
A formal series $\widehat f=\sum f_nq^n$ belongs to $\CM$ if
$|f_n|\le CA^nM_n$ for some constants. A function $f$ has a strong
$\M$-asymptotic expansion $\widehat f$ on $S_{\eta,r}$ if
\begin{equation}
 \left|f(q)-\sum_{n=0}^{N-1}f_nq^n\right|
 \le CA^NM_N|q|^N\quad(N\ge1,\ q\in S_{\eta,r}),
 \label{eq:strong}
\end{equation}
with constants independent of $N$ and $q$. For a sectorial germ this is
required on every proper left sector, allowing the radius and constants
to depend on that sector.
\end{definition}

\begin{lemma}[Elementary weight facts]
\label{lem:weightfacts}
For an admissible weight,
\begin{equation}
 M_aM_b\le M_{a+b},\qquad M_n\le\mu_N^n\quad(0\le n\le N).
 \label{eq:weightfacts}
\end{equation}
Multiplication by $q$ and, for zero-constant-term expansions, division by
$q$ preserve the coefficient and strong asymptotic classes. Strong
$\M$-asymptotics imply formal membership in $\CM$.
\end{lemma}
\begin{proof}
Multiply the ordered quotients to obtain the inequalities. Multiplication
by $q$ shifts an index down, and $M_{n-1}\le M_n$; division shifts up,
which is absorbed by the bound involving $E$. For coefficient necessity,
divide the remainder after degree $N-1$ by $q^N$ and approach zero along
an interior ray. The expansion one order further identifies the limit
as $f_N$, and gives $|f_N|\le CA^NM_N$.
\end{proof}

\begin{lemma}[Factorial--logarithmic weights]
\label{lem:logweights}
For $s>0$ and $\tau\in\R$, the weight \eqref{eq:logweight}, with a finite
initial quotient adjustment if needed, is admissible, and
\begin{align}
 \mu_n&\sim n^s(\log n)^\tau,\label{eq:muasymp}\\
 \log M_n&=sn\log n-sn+\tau n\log\log n+o(n).
 \label{eq:Masymp}
\end{align}
It is geometrically equivalent to
$(n!)^s(\log(\ee+n))^{\tau n}$.
\end{lemma}
\begin{proof}
The function $x^s(\log(\ee+x))^\tau$ is increasing and at least one
for all sufficiently large $x$. Keep these quotients from an index $J$
onward and choose nondecreasing quotients between $1$ and $\mu_J$
before that index. This changes the later weights by a constant factor.
Stirling's formula and
\[
 \sum_{\ell=1}^n\log\log(\ee+\ell)
 =n\log\log n+O(n/\log n)
\]
give \eqref{eq:Masymp}. The sum estimate follows by splitting at $\sqrt n$
and comparing with an integral, then integrating by parts.
Thus $n\log\mu_n-\log M_n=sn+o(n)$, which proves the $D$ bound.
The quotient $\mu_{n+1}$ has polynomial times logarithmic growth and is
bounded by $E^{n+1}$ after enlarging $E$. The same sum estimate proves
geometric equivalence with the simpler weight.
\end{proof}

\section{A weighted implicit remainder theorem}
\label{sec:implicit}
The uniform-in-order step uses only Cauchy estimates, contraction, and
an analytic majorant. No summability hypothesis on the kernel is used.

\begin{theorem}[Uniform weighted implicit expansion]
\label{thm:implicit}
Let $\M$ be admissible. Suppose $F$ is holomorphic on
$S_{\eta,r_0}\times D(1,\delta)$, with coefficient functions $a_n$
holomorphic in the disk, $a_0=1$, and uniform bounds
\begin{align}
 |a_n(v)|&\le CA^nM_n\quad(n\ge1),\label{eq:implicit-coeff}\\
 \left|F(q,v)-\sum_{n=0}^{N-1}a_n(v)q^n\right|
 &\le CA^NM_N|q|^N\quad(N\ge1).\label{eq:implicit-rem}
\end{align}
After shrinking the radius, the fixed point $V=F(q,V)$ near $1$ exists,
is holomorphic, and has a strong $\M$-expansion equal to the unique
formal solution. Constants can be obtained from the stated data.
\end{theorem}
\begin{proof}
Increase constants so that $A,C,D\ge1$.

\paragraph{Fixed-disk existence.}
The $N=1$ estimate gives $F(q,v)=1+O(q)$ uniformly. Cauchy's estimate
gives $\partial_vF=O(q)$ on $D(1,\delta/2)$. Shrinking the radius makes
$F$ preserve the closed smaller disk with Lipschitz constant at most
$1/2$. Iteration gives a unique fixed point, holomorphic by locally
uniform convergence.

\paragraph{Formal majorization.}
Write $V=1+w$. Cauchy's estimate gives
\[
 a_n(1+w)=\sum_{j\ge0}a_{n,j}w^j,
 \qquad |a_{n,j}|\le CA^nM_n\delta^{-j}.
\]
In total degree $N$, the weights of a term $q^nw_{i_1}\cdots w_{i_j}$
are bounded by $M_nM_{i_1}\cdots M_{i_j}\le M_N$. Induction therefore
dominates $|w_N|/M_N$ by $[z^N]H(z)$, where
\begin{equation}
 H=\frac{CAz}{1-Az}\frac1{1-H/\delta},\qquad H(0)=0.
 \label{eq:analyticmajorant}
\end{equation}
This analytic scalar equation is contractive on $|H|\le\delta/2$ for
$|z|\le\min\{(2A)^{-1},\delta/(16CA)\}$. Its coefficients are
nonnegative. Thus there are $K,B\ge1$ such that the unique formal solution
$\Vhat=\sum v_nq^n$ satisfies
\begin{equation}
 |v_n|\le KB^nM_n\quad(n\ge0).
 \label{eq:Vcoeff}
\end{equation}
Formal uniqueness follows because every nonconstant kernel term has
positive $q$-order.

\paragraph{Shrinking analytic disks.}
Put $F_{<N}=\sum_{n<N}a_n(v)q^n$ and
$\rho_N=c/(A\mu_N)$, where
$0<c\le\min\{1/4,\delta/(16C)\}$. On $|q|\le\rho_N$,
\[
 \sum_{n=1}^{N-1}CA^nM_n|q|^n\le C\sum_{n\ge1}c^n
 \le Cc/(1-c).
\]
Cauchy's estimate in $v$ makes $F_{<N}$ contractive on
$\overline D(1,\delta/2)$, decreasing $c$ if needed. Its fixed point
$V_N$ is analytic in $D(0,\rho_N)$ and has the same Taylor coefficients
as $\Vhat$ through degree $N-1$.
For sectorial $q$ with $|q|\le\rho_N/2$, subtract the fixed-point
equations to get
\begin{equation}
 |V-V_N|\le2CA^NM_N|q|^N.
 \label{eq:compareVN}
\end{equation}
Cauchy's estimate for $V_N-1$, bounded by $\delta/2$, gives
\begin{equation}
 \left|V_N-\sum_{n<N}v_nq^n\right|
 \le\delta(|q|/\rho_N)^N
 \le\delta(AD/c)^NM_N|q|^N.
 \label{eq:VNtail}
\end{equation}
This proves the estimate in the small-$q$ range associated with $N$.

\paragraph{The complementary range.}
When $|q|\ge\rho_N/2$, for $0\le n<N$,
\begin{equation}
 \frac{M_n|q|^n}{M_N|q|^N}
 =\frac{M_n}{\mu_N^n}\frac{\mu_N^N}{M_N}
      (|q|\mu_N)^{-(N-n)}
 \le D^N(2A/c)^{N-n}.
 \label{eq:largerange}
\end{equation}
Combine this with \eqref{eq:Vcoeff} and the boundedness of $V$.
Summing and using $N+1\le2^N$ bounds the normalized remainder by a
constant times $[2D\max\{B,2A/c,1\}]^N$. Together with
\eqref{eq:compareVN}--\eqref{eq:VNtail}, this is a single all-orders
strong remainder estimate on the smaller sector.
\end{proof}

\begin{remark}
This proof permits a geometric enlargement of the weight. It preserves
class membership, not necessarily the optimal numerical coefficient or
remainder type. The shrinking-disk Cauchy estimate is one explicit source
of such a loss of constants.
\end{remark}

\section{The desingularized countable kernel}
\label{sec:kernel}
Fix $0<\eta<1$, choose $0<\delta<\eta$, and put
$c=\eta-\delta>0$, $B=1+\delta$. For $q\in S_{\eta,r}$ and
$|v-1|<\delta$,
\begin{equation}
 \Rea(qv)\le-c|q|.
 \label{eq:negativeproduct}
\end{equation}
All the following estimates are uniform on this product domain.

\begin{proposition}[Sectorial existence with arbitrary slopes]
\label{prop:existence}
For every nonnegative slope sequence, $F$ in \eqref{eq:desing} is
holomorphic on $S_{\eta,r}\times D(1,\delta)$ when $r<1$, and
\begin{align}
 |F(q,v)-1|&\le\lambda_1B|q|+\frac{|q|}{1-|q|},\label{eq:Fclose}\\
 |\partial_vF(q,v)|&\le\lambda_1|q|
                  +\frac{|q|}{\ee c(1-|q|)}.\label{eq:Fderivative}
\end{align}
For a sufficiently small radius it has a unique fixed point near $1$.
The corresponding $U=qV$ solves \eqref{eq:model}, has $U=q+O(q^2)$,
and is unique among solutions with that behavior on a proper left sector.
Constructions on different sectors agree near zero on their overlaps.
\end{proposition}
\begin{proof}
The $j$th term is bounded by $|q|^{j-1}$, giving normal convergence.
For $j=1$ use $e^z-1=z\int_0^1e^{tz}\dd t$ and
$\Rea z\le0$; the remaining terms have a geometric majorant.
For the derivative handle $j=1$ directly. For $j\ge2$, the inequality
$t e^{-ct|q|}\le1/(\ee c|q|)$, $t\ge0$, bounds the derivative term
by $|q|^{j-1}/(\ee c)$. Summing proves the two estimates.
They tend to zero, giving contraction on a fixed closed disk about $1$.
The first estimate then gives $V=1+O(q)$. Any other solution with the
stated behavior has $U/q$ in the same disk near zero, proving uniqueness
and compatibility.
\end{proof}

\begin{proposition}[Exact uniform kernel remainder]
\label{prop:kernelremainder}
Define
\begin{equation}
 a_n(v)=\sum_{j=1}^{n+1}
        \frac{(\lambda_jv)^{n+1-j}}{(n+1-j)!},\qquad n\ge0.
 \label{eq:an}
\end{equation}
Then $a_0=1$, $|a_n(v)|\le a_n(B)$, and for every $N\ge1$,
\begin{equation}
 \boxed{\quad
 \left|F(q,v)-\sum_{n=0}^{N-1}a_n(v)q^n\right|
 \le\frac{a_N(B)}{1-|q|}|q|^N.
 \quad}
 \label{eq:exactkernelremainder}
\end{equation}
No growth restriction on the slopes is needed.
\end{proposition}
\begin{proof}
For $\Rea z\le0$, $m\ge1$, the integral Taylor remainder gives
\begin{equation}
 \left|e^z-\sum_{h=0}^{m-1}\frac{z^h}{h!}\right|
 =\left|\frac{z^m}{(m-1)!}
       \int_0^1(1-t)^{m-1}e^{tz}\dd t\right|
 \le\frac{|z|^m}{m!}.
 \label{eq:exptail}
\end{equation}
For $j\le N$, apply this with $z=\lambda_jqv$ and $m=N+1-j$,
then multiply by $q^{j-1}$. The combined bound is
\[
 |q|^N\sum_{j=1}^N\frac{(B\lambda_j)^{N+1-j}}{(N+1-j)!}.
\]
The terms $j>N$ contribute at most $|q|^N/(1-|q|)$.
The missing $j=N+1$ summand of $a_N(B)$ equals $1$, so their sum is
bounded by \eqref{eq:exactkernelremainder}.
\end{proof}

The estimate holds on a fixed disk of dependent-variable values, not
only along the unknown solution. This is the uniformity required in
Theorem~\ref{thm:implicit}. Substitution of a formal polynomial into the
original kernel would not establish it.

\section{Positivity and the coefficient-to-remainder principle}
\label{sec:bridge}
The following algebraic comparison does not depend on the sharper
coefficient asymptotics in the predecessor drafts.

\begin{lemma}[Formal coefficients and a positive subfamily]
\label{lem:positivity}
Equation~\eqref{eq:model} has a unique solution in $q\R[[q]]$.
It has $u_1=1$, $u_n\ge1$, and
\begin{equation}
 u_n=\sum_{\substack{\boldsymbol m\text{ finite}\\\sum j m_j=n}}
 \frac{(\sum_j\lambda_jm_j)^{|\boldsymbol m|-1}}{\prod_jm_j!},
 \qquad |\boldsymbol m|=\sum_jm_j,
 \label{eq:Lagrange}
\end{equation}
where $0^0=1$. In particular,
\begin{equation}
 a_n(1)\le u_{n+1}\quad(n\ge0),
 \label{eq:positivesubfamily}
\end{equation}
and for $j\ge2$, $h\ge0$,
\begin{equation}
 u_{j+h}\ge\frac{\lambda_j^h}{h!}.
 \label{eq:oneactionlower}
\end{equation}
\end{lemma}
\begin{proof}
The defining map raises the order of a difference of formal candidate
solutions by at least one, giving existence and uniqueness by iteration.
Introduce a parameter $t$ and solve $U=t\Phi(q,U)$, where
$\Phi(q,z)=\sum_{j\ge1}q^je^{\lambda_jz}$. Lagrange inversion gives
\[
 U=\sum_{k\ge1}\frac{t^k}{k}[z^{k-1}]\Phi(q,z)^k.
\]
At each $q$-degree all sums are finite. Expanding by multiplicities
cancels the multinomial and derivative factorials, yielding
\eqref{eq:Lagrange} when $t=1$.

All terms are nonnegative. One action $j$ and $h$ actions $1$ contribute
$(\lambda_j+h\lambda_1)^h/h!$, proving \eqref{eq:oneactionlower}.
For \eqref{eq:positivesubfamily}, put $N=n+1$. The configurations with
one action $j\ge2$ and $N-j$ actions $1$ are distinct as $j$ varies.
They dominate the corresponding terms of $a_{N-1}(1)$. The all-action-$1$
configuration contributes $(N\lambda_1)^{N-1}/N!$, at least
$\lambda_1^{N-1}/(N-1)!$ for $N\ge2$, since their ratio when nonzero
is $N^{N-2}$. The case $N=1$ is equality. The single action $n$ also
proves $u_n\ge1$.
\end{proof}

\begin{theorem}[Coefficient-to-remainder equivalence]
\label{thm:bridge}
Let $\M$ be admissible and $U$ the canonical left-sector solution.
The following statements are equivalent:
\begin{enumerate}[label=(\roman*)]
\item $\Uhat\in\CM$;
\item $a_n(1)\le CA^nM_n$ for some $C,A>0$;
\item $U$ has strong $\M$-asymptotics $\Uhat$ on every proper left sector;
\item $U$ has that strong expansion on at least one proper left sector.
\end{enumerate}
\end{theorem}
\begin{proof}
By \eqref{eq:positivesubfamily} and the shift bound for $M_{n+1}$,
(i) implies (ii). Since $a_n(B)\le B^na_n(1)$ for $B\ge1$,
Proposition~\ref{prop:kernelremainder} and Theorem~\ref{thm:implicit}
turn (ii) into strong $\M$-asymptotics of $V=U/q$ on every proper left
sector. Multiplication by $q$ preserves the class, and formal uniqueness
identifies the coefficients; hence (iii). The implication (iii) to (iv)
is immediate, and (iv) to (i) is coefficient necessity from
Lemma~\ref{lem:weightfacts}.
\end{proof}

The theorem proves that the canonical realization introduces no loss of
weight class beyond geometric constants. Conversely, failure of formal
membership rules out the corresponding strong expansion on every proper
left sector. Positivity is essential in the present proof: it prevents
cancellation from hiding a large kernel coefficient.

\section{Determining the factorial--logarithmic scale}
\label{sec:classification}
We prove the class estimates independently, without invoking an exact
numerical Gevrey type or a multiplicative saddle equivalent from the repo.

\begin{lemma}[Growth of the kernel coefficients]
\label{lem:kernelgrowth}
Fix $p>1$, $\beta\in\R$. If
$\lambda_j\le Cj^p(\log(\ee+j))^\beta$ eventually, then for each
fixed $B\ge1$ there are $K,A>0$ such that
\begin{equation}
 a_n(B)\le KA^nM_{p-1,\beta-p;n}.
 \label{eq:kernelgrowthupper}
\end{equation}
Under the two-sided estimate \eqref{eq:twosidedslopes},
\begin{equation}
 \log a_n(1)=(p-1)n\log n+(\beta-p)n\log\log n+O(n).
 \label{eq:kernelgrowthlog}
\end{equation}
\end{lemma}
\begin{proof}
Enlarge a constant to cover all exceptional slopes. Put $L=\log n$ and
consider the term with index $j$ and exponent $h=n+1-j$. The term $h=0$
is $1$. For $h\ge1$, using $\log h!\ge h\log h-h$ gives
\begin{equation}
 \log\frac{(B\lambda_j)^h}{h!}
 \le h\{C_0+p\log j+\beta\log\log(\ee+j)-\log h\}.
 \label{eq:optimization-start}
\end{equation}
We bound the maximum uniformly in $j$.

For $j<\sqrt n$, the right side is at most
$(p/2-1)nL+O(n\log L)$, smaller than the claimed bound by a positive
multiple of $nL$. For $j>n/2$, $h\le n/2+1$ and
$\log\lambda_j\le pL+O(\log L)$. The function
$h(pL+O(\log L)-\log h+1)$ increases on this range for large $n$,
so its maximum is at most
$\tfrac12(p-1)nL+O(n\log L)$, also smaller than required.

In the middle range $\sqrt n\le j\le n/2$,
$\log\log(\ee+j)=\log L+O(1)$ and $\log h=L+O(1)$.
Expanding $h=n+1-j$, with $s=p-1$, gives
\begin{align*}
 \log\frac{(B\lambda_j)^h}{h!}
 \le{}&snL+\beta n\log L+pn\log(j/n)-sjL\\
      &+|\beta|j\log L-pj\log(j/n)+O(n)\\
 \le{}&snL+\beta n\log L+pn\log(j/n)-\tfrac s2jL+O(n).
\end{align*}
We used $-j\log(j/n)\le n/\ee$ and
$|\beta|\log L\le sL/2$ eventually. Setting $t=jL/n$ turns the last
line into
\[
 snL+(\beta-p)n\log L+n\{p\log t-st/2\}+O(n).
\]
The expression in braces is bounded above for $t>0$. Summing $n+1$
terms changes the logarithm by at most $\log(n+1)$, and
Lemma~\ref{lem:logweights} proves the upper bound.

For a lower bound, take $j=\lfloor n/\log n\rfloor$ and $h=n+1-j$.
The two-sided slope assumption gives
\[
 \log\lambda_j=p\log n+(\beta-p)\log\log n+O(1),
 \qquad h=n+O(n/\log n).
\]
Stirling's formula for $h!$ yields the matching lower bound in
\eqref{eq:kernelgrowthlog}. Zero slopes in the upper-bound argument are
simply omitted.
\end{proof}

\begin{theorem}[Exact Gevrey class criterion]
\label{thm:gevrey}
For every $s>0$,
\begin{equation}
 \Uhat\in\mathcal C_{(n!)^s}
 \quad\Longleftrightarrow\quad
 \lambda_j=O((j\log(j+1))^{s+1}).
 \label{eq:gevreycriterion}
\end{equation}
Both conditions are equivalent to strong Gevrey-$s$ asymptotics of the
canonical left-sector solution.
\end{theorem}
\begin{proof}
For sufficiency use Lemma~\ref{lem:kernelgrowth} with $p=s+1$,
$\beta=p$, giving logarithmic exponent zero, and then
Theorem~\ref{thm:bridge}.
For necessity suppose $u_n\le CA^n(n!)^s$. For each large $j\ge2$ set
$h=\lceil j\log j\rceil$ and $n=j+h$. By \eqref{eq:oneactionlower},
\[
 \lambda_j^h\le h!CA^n(n!)^s,
 \qquad
 \lambda_j\le C^{1/h}A^{n/h}h\,n^{sn/h}.
\]
Now $n/h=1+O(1/\log j)$, $n\asymp j\log j$, and
$n^{s(n/h-1)}=O(1)$. Thus
$\lambda_j\le C'(j\log j)^{s+1}$.
\end{proof}

\begin{proposition}[The analytic endpoint]
\label{prop:analytic}
The formal solution has positive convergence radius exactly when
$\lambda_j=O(j)$.
\end{proposition}
\begin{proof}
If $\lambda_j\le Cj$, the original kernel converges normally in a small
bidisk, its terms bounded by $(|q|e^{C|U|})^j$. The analytic implicit-
function theorem applies at $(0,0)$. Conversely an exponential bound
$u_n\le CA^n$ and \eqref{eq:oneactionlower} with $h=j$ give
$\lambda_j\le C^{1/j}A^2(j!)^{1/j}=O(j)$.
\end{proof}
The positive-$s$ criterion must not be extrapolated to $s=0$.
For $\lambda_j=j\log(j+1)$ the series is divergent but belongs to
every positive Gevrey class, and its canonical realization has each
corresponding strong expansion.

\begin{theorem}[Optimality within the logarithmic scale]
\label{thm:logoptimal}
Under \eqref{eq:twosidedslopes}, with $s=p-1$, $\tau=\beta-p$,
\begin{equation}
 \log u_n=sn\log n+\tau n\log\log n+O(n).
 \label{eq:ulog}
\end{equation}
The canonical solution has strong $\M_{s,\tau}$-asymptotics, but no
strong $\M_{s,\tau-\varepsilon}$-asymptotics for $\varepsilon>0$ on
any proper left sector. Decreasing $s$ cannot be compensated by any fixed
logarithmic exponent.
\end{theorem}
\begin{proof}
The upper bound follows from Lemma~\ref{lem:kernelgrowth} and the formal
majorant in Theorem~\ref{thm:implicit}; the lower bound follows from
$a_{n-1}(1)\le u_n$ and \eqref{eq:kernelgrowthlog}. Shifting the index
by one changes the expression by $o(n)$. Strong membership follows from
Theorem~\ref{thm:bridge}.
Reducing $\tau$ by $\varepsilon$ subtracts
$\varepsilon n\log\log n$ from the logarithm of the weight, which no
geometric factor can absorb. Reducing $s$ creates a larger deficit of
order $n\log n$. Coefficient necessity rules out both improvements.
\end{proof}

\section{Actual optimal-truncation bounds}
\label{sec:optimal}
Define the associated function
\begin{equation}
 h_{\M}(t)=\inf_{N\ge0}M_Nt^N,\qquad t>0.
 \label{eq:hM}
\end{equation}
The all-orders estimates now allow the truncation index to vary with $q$.

\begin{proposition}[Associated error scale]
\label{prop:hM}
For $\M_{s,\tau}$ let $n(t)$ be an integer minimizing $M_Nt^N$.
As $t\downarrow0$,
\begin{align}
 n(t)&\sim s^{\tau/s}t^{-1/s}(\log(1/t))^{-\tau/s},\label{eq:nt}\\
 \log h_{\M}(t)&\sim-sn(t).\label{eq:hMasymp}
\end{align}
\end{proposition}
\begin{proof}
Consecutive terms have ratio $t\mu_{N+1}$; a minimum occurs at the
crossing of one. Since adjacent quotient ratios tend to one,
$t n(t)^s(\log n(t))^\tau\to1$. Substituting the right side of
\eqref{eq:nt} verifies this inverse asymptotic. Its logarithm has leading
term $\log(1/t)/s$. At that index, \eqref{eq:Masymp} gives
\[
 \log(M_{n(t)}t^{n(t)})
 =n(t)\{s\log n(t)+\tau\log\log n(t)+\log t-s\}+o(n(t))
 =-sn(t)+o(n(t)).
\]
A one-index tie does not affect the equivalent.
\end{proof}

\begin{corollary}[Realized small errors]
\label{cor:actualerror}
If $U$ has a strong $\M_{s,\tau}$ bound with geometric constant $A$,
choose $N=n(A|q|)$. Then
\begin{equation}
 \left|U(q)-\sum_{n<N}u_nq^n\right|
 \le C\exp[-c|q|^{-1/s}(\log(1/|q|))^{-\tau/s}]
 \label{eq:actualerror}
\end{equation}
for sufficiently small $q$ on the chosen proper sector.
\end{corollary}
\begin{proof}
Take the infimum of the strong all-orders estimate and apply
Proposition~\ref{prop:hM}. The geometric constant changes the positive
constant in the exponential, not its power or logarithmic correction.
\end{proof}

For quadratic slopes $\lambda_j=j^2$, $s=1$, $\tau=-2$, so
\begin{equation}
 \left|U(q)-\sum_{n<N(q)}u_nq^n\right|
 \le C\exp\left[-c\frac{(\log(1/|q|))^2}{|q|}\right],
 \qquad N(q)\asymp\frac{(\log(1/|q|))^2}{|q|}.
 \label{eq:quadraticactual}
\end{equation}
This is an actual error estimate, not merely a formal least-term law.
It implies strong Gevrey-$1$ bounds with arbitrarily small positive
geometric type, since
$A^N(\log(N+\ee))^{-2N}\le C_\varepsilon\varepsilon^N$ for every
$\varepsilon>0$. The Taylor series is nevertheless divergent.

Under $q=e^{-x}$ the formal solution is
$\sum_{n\ge1}u_ne^{-nx}$. Proper left sectors correspond to horizontal
strips about odd multiples of $\pi$ in $\Im x$. With $X=\Rea x$,
\eqref{eq:actualerror} becomes
\begin{equation}
 O(\exp[-c e^{X/s}X^{-\tau/s}]),\qquad
 N(X)\asymp e^{X/s}X^{-\tau/s}.
 \label{eq:transserieserror}
\end{equation}
These are complex-strip realizations, not claims on the real-$x$ axis.

\section{Sharp negative-direction summability}
\label{sec:summability}
We use one classical external analytic characterization, with its
hypotheses verified by the preceding proofs.

\paragraph{Classical sectorial characterization.}
For $s=1/k>0$, a formal series is $k$-summable in direction $d$ exactly
when it is the strong Gevrey-$s$ expansion of a holomorphic function on a
sectorial region centered at $d$ with opening strictly greater than
$\pi s$. The function is the unique sum there. Equivalently its
moment-normalized Borel transform
\begin{equation}
 \mathcal B_k\widehat f(\xi)
 =\sum_{n\ge0}\frac{f_n\xi^n}{\Gamma(1+n/k)}
 \label{eq:kBorel}
\end{equation}
has the required analytic continuation and exponential-order-$k$ growth
near direction $d$, and generalized Laplace reconstruction recovers the
function. We use the local-sector formulation in Lastra--Malek--Sanz,
Definition 3.3, Theorem 3.18, and Remark 3.19(i) \cite{lms}. The weights
$(n!)^{1/k}$ and $\Gamma(1+n/k)$ are geometrically equivalent.

\begin{theorem}[Exact summability criterion for $k>1$]
\label{thm:ksum}
For every $k>1$,
\begin{equation}
 \boxed{\quad
 \Uhat\text{ is $k$-summable in direction }\pi
 \quad\Longleftrightarrow\quad
 \lambda_j=O((j\log(j+1))^{1+1/k}).
 \quad}
 \label{eq:ksumcriterion}
\end{equation}
Whenever these conditions hold, its sum is the canonical solution $U$.
\end{theorem}
\begin{proof}
Set $s=1/k<1$. The slope condition and Theorem~\ref{thm:gevrey} give
strong Gevrey-$s$ asymptotics on every proper left sector. Choose $\eta$
so small that $2\arccos\eta>\pi s$. The classical characterization
applies there and identifies the sum with $U$.
Conversely $k$-summability implies a Gevrey-$1/k$ coefficient bound.
Necessity in Theorem~\ref{thm:gevrey} gives the slope condition.
\end{proof}

\begin{corollary}[A logarithmically refined sum]
\label{cor:Msum}
Under \eqref{eq:twosidedslopes} with $1<p<2$, the formal solution is
$\M_{p-1,\beta-p}$-summable in direction $\pi$ in the standard
generalized ultraholomorphic sense. Its generalized sum is $U$.
\end{corollary}
\begin{proof}
The factorial--logarithmic weights, with equivalent initial normalization,
are strongly regular and have opening index $s=p-1$; these weight facts
and their generalized reconstruction theory are established in
\cite{lms,sanz}. Theorem~\ref{thm:logoptimal} supplies strong asymptotics
for this weight. Since $s<1$, an available left sector has opening
greater than $\pi s$, so the characterization applies.
\end{proof}

For $\lambda_j=j^p$, $1<p<2$, the critical ordinary summation order is
$k=1/(p-1)$. For slopes comparable to
$j^p(\log(\ee+j))^\beta$, ordinary summability at that critical order
holds exactly when $\beta\le p$. When $\beta>p$ it fails, but the
logarithmically refined sum still exists. Every ordinary order satisfying
$p-1<1/k<1$ remains available.

\subsection{The critical aperture is not resolved by this argument}
For $s=1$ the directly available sectorial region has opening $\pi$, not
strictly greater. Strong Gevrey-$1$ estimates on its proper subsectors
therefore do not establish ordinary Borel summability by this theorem.
The quadratic type-zero improvement does not give uniform control up to
the imaginary boundary rays.

A Nevanlinna--Sokal criterion on a suitable tangent-disk domain could be
relevant, but its geometric and uniformity hypotheses have not been
established here \cite{sokal}. Thus ordinary negative-direction Borel
summability for $\lambda_j=j^2$ remains a stated research question, not
an implicit corollary.
% ed. (2026-09-29): the stated research question is settled by two sibling packages.
\begin{ednote}
This question is settled for $\lambda_j=j^2$ by two other packages of
the series. In the fine (half-strip) sense the answer is positive: \nolinkurl{Negative_Ray_Summation_Exponential_Feedback/}~\cite{ed:negativeray}, Theorem~\texttt{thm:threshold}, available at this article's snapshot. In
the ordinary (angular) sense it is negative: \nolinkurl{Natural_Boundaries_Quadratic_Exponential_Feedback/}~\cite{ed:naturalboundaries}, Theorem~\texttt{thm:main},
proves that the literal inverse has a natural boundary on the imaginary
diameter and that neither $\Uhat$ nor its inverse is $1$-summable in any
open arc of directions containing $\pi$. Uniform estimates for $U$ up to
the boundary rays are not provided by either.
\end{ednote}

\subsection{The literal positive-real obstruction}
\begin{proposition}
\label{prop:positive}
If $\sup_j\lambda_j/j=\infty$, the literal infinite sum
\eqref{eq:model} has no solution for all sufficiently small positive
real $q$ with $U(q)/q\to1$.
\end{proposition}
\begin{proof}
Such a solution would have $\Rea U(q)>0$ for sufficiently small positive
$q$. Fix one such $q$ and choose a subsequence with
$\lambda_j/j\to\infty$. Then
\[
 \log|q^je^{\lambda_jU(q)}|=j\log q+\lambda_j\Rea U(q)
 \longrightarrow+\infty
\]
along that subsequence, so the summands do not tend to zero.
\end{proof}
This concerns convergence of the defining kernel. Analytic continuation
of a function constructed elsewhere, or a generalized sum in another
direction, is a different question. We do not infer failure of every
possible continuation from this test.

\section{Sectorial compositional inversion}
\label{sec:inverse}
The inverse result gives an actual inverse germ with strong remainders,
not merely a transfer of a formal regularity label.

\begin{lemma}[Formal weighted reversion]
\label{lem:formalreversion}
For an admissible weight and
$\widehat f(q)=q+\sum_{n\ge2}f_nq^n$,
\[
 \widehat f\in\CM\quad\Longleftrightarrow\quad
 \widehat f^{\langle-1\rangle}\in\CM.
\]
\end{lemma}
\begin{proof}
Write $\widehat f(q)=q(1+\widehat h(q))$ and its proposed inverse as
$q=zW(z)$, $W(0)=1$. The inverse equation is
\begin{equation}
 W=1-W\widehat h(zW).
 \label{eq:inverse-formal}
\end{equation}
The coefficient of each positive power of $z$ is a polynomial in $W$.
On a fixed disk about $1$ it is bounded by a geometric factor times the
corresponding coefficient of $\widehat h$. Division by $q$ preserves the
weight class, so the formal-majorant step of
Theorem~\ref{thm:implicit} proves $W\in\CM$. Multiplication by $z$
proves one implication; applying it again to the inverse proves the other.
\end{proof}

\begin{theorem}[Strong inverse realization]
\label{thm:inverse}
There is a canonical sectorial inverse $Q$ on a left sectorial germ,
with $Q(z)=z+O(z^2)$. For every admissible weight,
\begin{equation}
 \Uhat\in\CM\quad\Longleftrightarrow\quad
 \Qhat\in\CM\quad\Longleftrightarrow\quad
 Q\text{ has strong $\M$-asymptotics}.
 \label{eq:inverseequiv}
\end{equation}
The Gevrey, summability, and optimal logarithmic-scale conclusions all
hold for the inverse. When directional sums are asserted, they are actual
compositional inverses.
\end{theorem}
\begin{proof}
Choose nested proper left sectors. On the larger one $U(q)-q=O(q^2)$;
Cauchy's estimate on disks of radius proportional to $|q|$ gives
$U'(q)-1=O(q)$ on the smaller one. For $z$ in a still smaller sector,
solve $q=z-(U(q)-q)$ in a disk centered at $z$ of radius $K|z|^2$.
For fixed sufficiently large $K$ and small $z$, the disk stays in the
preceding sector, the map preserves it, and its derivative has modulus
less than $1/2$. This constructs a holomorphic $Q=z+O(z^2)$ satisfying
$U(Q(z))=z$. Local inverse uniqueness gives the other composition identity
on compatible domains and the compatibility of sectorial restrictions.

Suppose now $\Uhat\in\CM$. Theorem~\ref{thm:bridge} gives a strong
$\M$-expansion for $V=U/q$. Choose a sufficiently small fixed disk
$D(1,\delta)$ so that $zv$ stays in the larger sector for $z$ in the
chosen smaller sector and $v$ in that disk. For $Q=zW$, the equation is
\begin{equation}
 W=1-W(V(zW)-1).
 \label{eq:inverse-analytic}
\end{equation}
If $V(q)\sim\sum v_nq^n$, the right side has coefficient polynomials
$-v_nW^{n+1}$ for $n\ge1$. These satisfy the required weighted estimates
on the fixed disk. Its remainder is bounded by
$CA^NM_N|zW|^N|W|$, hence by the same weight with a geometric enlargement.
Theorem~\ref{thm:implicit} yields strong asymptotics for $W$, and then
for $Q$.

Strong asymptotics imply $\Qhat\in\CM$, and formal weighted reversion
returns $\Uhat\in\CM$. A smaller weight for $Q$ would therefore give a
smaller weight for $U$, contrary to Theorem~\ref{thm:logoptimal}.
The summability statement follows by the same sector-opening argument
and uniqueness of the corresponding sum.
\end{proof}

\begin{remark}[Cancellation is not ignored]
We have proved weight-class optimality for the inverse, not a pointwise
asymptotic formula for its signed coefficients. Equation \eqref{eq:ulog}
is proved for the positive coefficients of $\Uhat$ and is not silently
transferred to $\Qhat$. Equality of their sharp numerical Gevrey types
is likewise not asserted.
\end{remark}

\section{Finite-action evaluation and executable checks}
\label{sec:numerics}
\subsection{A uniform analytic certificate}
Define
\begin{equation}
 F_L(q,v)=\sum_{j=1}^Lq^{j-1}e^{\lambda_jqv},\qquad V_L=F_L(q,V_L).
 \label{eq:finitekernel}
\end{equation}
The same disk and derivative bounds hold. One may choose
$\delta=\eta/4$, $c=3\eta/4$, $B=1+\delta$, and a radius such that
\begin{equation}
 \lambda_1Br+\frac r{1-r}\le\delta/2,
 \qquad \kappa(r):=\lambda_1r+\frac r{\ee c(1-r)}<1.
 \label{eq:numericconditions}
\end{equation}
Both finite and infinite maps preserve the chosen disk and have
Lipschitz constant at most $\kappa(r)$.

\begin{proposition}[Action and iteration error bounds]
\label{prop:certificate}
Under \eqref{eq:numericconditions}, with $|q|=r$,
\begin{equation}
 |U(q)-qV_L(q)|\le\frac{r^{L+1}}{(1-\kappa(r))(1-r)}.
 \label{eq:actioncertificate}
\end{equation}
If $v_0=1$, $v_{m+1}=F_L(q,v_m)$, then
\begin{equation}
 \boxed{\quad
 |U(q)-qv_m|\le
 \frac{r\kappa(r)^m}{1-\kappa(r)}|v_1-v_0|
 +\frac{r^{L+1}}{(1-\kappa(r))(1-r)}.
 \quad}
 \label{eq:fullcertificate}
\end{equation}
The action-tail bound and $\kappa$ are independent of the slopes
$\lambda_j$ for $j\ge2$. In the iteration bound, replacing the computed
increment $|v_1-v_0|$ by $\lambda_1Br+r/(1-r)$ makes the entire
right side independent of those slopes.
\end{proposition}
\begin{proof}
The kernel tail is at most $r^L/(1-r)$. Subtract the fixed-point equations,
use contraction, and multiply by $r$ to prove
\eqref{eq:actioncertificate}. Summing the geometric bound on successive
iteration increments gives
$|V_L-v_m|\le\kappa^m|v_1-v_0|/(1-\kappa)$; add the estimates.
\end{proof}

Finite-action truncation and Taylor truncation are different. The former
converges geometrically at fixed $q$ even when the Taylor series diverges;
the latter has the optimally chosen error bounds proved above.
For end-to-end certified numerics the finite exponential evaluations and
the right side of \eqref{eq:fullcertificate} must also be enclosed with
rounding errors. The supplied high-precision diagnostics do not use
interval arithmetic. Their reported error bounds are floating-point
evaluations of the proved analytic certificate.

\subsection{An integer recurrence}
For integer slopes let $b_n=n!u_n$ and
$E_{j,n}=n![q^n]e^{\lambda_j\Uhat(q)}$, with $E_{j,0}=1$.
Formal differentiation gives
\begin{align}
 b_n&=\sum_{j=1}^n\frac{n!}{(n-j)!}E_{j,n-j},\label{eq:integerU}\\
 E_{j,n}&=\lambda_j\sum_{m=1}^n
              \binom{n-1}{m-1}b_mE_{j,n-m}.\label{eq:integerexp}
\end{align}
At stage $n$ the first line uses previously computed exponential
coefficients; the second is then evaluated for every needed $j$.
All arithmetic is integral. Independent partition enumeration checks
\eqref{eq:Lagrange} at small orders.

\subsection{Recorded verification}
The bundled program was executed through degree $120$ for each of
$\lambda_j=0,j,j^2,j^3$. Its exact checks were:
\begin{center}
\begin{tabular}{@{}lr@{}}\toprule
Check & Number passed\\\midrule
Independent partition-formula comparisons & 48\\
Positive-subfamily inequalities $a_n(1)\le u_{n+1}$ & 480\\
Independent zero- and linear-slope benchmarks & 240\\\midrule
Total exact comparisons & 768\\\bottomrule
\end{tabular}
\end{center}
All passed. The linear benchmark uses the separate identity
$U=q(1+U)e^U$ and its one-variable Lagrange coefficient formula.
Full exact integer-normalized coefficients are included in CSV files.
For the quadratic model,
\begin{align*}
 \Uhat(q)={}&q+2q^2+\tfrac{15}{2}q^3+\tfrac{107}{3}q^4
 +\tfrac{1555}{8}q^5+\tfrac{17362}{15}q^6\\
 &+\tfrac{5288311}{720}q^7+\tfrac{15414767}{315}q^8+\cdots;
\end{align*}
for the cubic model the first terms are
\[
 q+2q^2+\tfrac{23}{2}q^3+\tfrac{269}{3}q^4
        +\tfrac{21433}{24}q^5+\tfrac{53144}{5}q^6+\cdots.
\]

At $100$ decimal digits the program also checked the kernel remainder
inequality in $48$ cases and compared six finite-action solves with
longer, more accurate solves. The reported approximations use $L=18$,
$m=24$. A sample is
\begin{center}\small
\begin{tabular}{@{}clll@{}}\toprule
$p$ & $q$ & Approximate $U(q)$ & Analytic bound\\\midrule
2 & $-0.01$ & $-0.0098071616822383391168465501$ & $1.031\cdot10^{-38}$\\
2 & $-0.005$ & $-0.0049509157982251340457345151$ & $1.937\cdot10^{-44}$\\
3 & $-0.01$ & $-0.0098106832787087467452630685$ & $1.031\cdot10^{-38}$\\
3 & $-0.005$ & $-0.0049513840936461045663286861$ & $1.937\cdot10^{-44}$\\
\bottomrule\end{tabular}
\end{center}
The bound concerns the internally retained high-precision approximation,
not the shortened printed decimal. CSV files record additional digits,
complex tests, and observed comparison errors. Numerical agreement does
not prove an asymptotic or summability theorem.

\section{Examples and boundary cases}
\label{sec:examples}
For $\lambda_j=j^{3/2}$ the optimal weight has $s=1/2$, $\tau=-3/2$.
The series is $2$-summable in direction $\pi$, with sum equal to the
literal left-sector fixed point. Optimized Taylor truncation gives
\[
 O(\exp[-c|q|^{-2}(\log(1/|q|))^3]).
\]
The logarithmically refined sum records the sharper natural weight.

Fix $k>1$. For slopes $(j\log(j+1))^{1+1/k}$, negative-direction
$k$-summability holds. Multiplication by $(\log(j+1))^\varepsilon$,
$\varepsilon>0$, destroys it by the necessary slope criterion.
The sectorial solution still exists, and a logarithmically refined
summation method or slightly smaller ordinary summation order remains
available.

For $\lambda_j=j^2$, strong Gevrey-$1$ asymptotics and
\eqref{eq:quadraticactual} are proved, but ordinary negative-direction
Borel summability is not decided here. For $\lambda_j=j^3$, the
natural weight is $M_{2,-3}$, giving
\[
 O(\exp[-c|q|^{-1/2}(\log(1/|q|))^{3/2}]).
\]
The directly available sector opening does not suffice to deduce
summability at the critical order. No failure of summability is inferred
merely from that limitation of the argument.

For $\lambda_j=e^{\sqrt j}$, the slope condition fails for every finite
$s>0$. No proper left sector therefore has a strong finite-order Gevrey
expansion of the canonical solution. The sectorial existence theorem and
finite-action certificate nevertheless remain valid. Other weight
classes may apply, but their hypotheses must be checked.

\begin{center}\small
\begin{tabular}{@{}p{0.24\textwidth}p{0.34\textwidth}p{0.34\textwidth}@{}}
\toprule
Slopes & Strong remainder information & Summability proved here\\\midrule
$O(j)$ & Convergent Taylor series & Analytic germ\\
$j\log(j+1)$ & Every positive Gevrey class & Every $k>1$, direction $\pi$\\
$j^p$, $1<p<2$ & $M_{p-1,-p}$, optimal in scale & Critical $k=1/(p-1)$ and refined sum\\
$j^2$ & $M_{1,-2}$ and \eqref{eq:quadraticactual} & Critical case left open\\
$j^3$ & $M_{2,-3}$, optimal in scale & No critical-order conclusion\\
$e^{\sqrt j}$ & No finite Gevrey order & No finite-order claim\\\bottomrule
\end{tabular}
\end{center}
After the first row all formal series are divergent. Every row still
has the canonical left-sector solution.

\section{Proof status and further research}
\label{sec:future}
\subsection{Statement-level proof ledger}
The sectorial construction, kernel remainder estimate, weighted implicit
theorem, positive coefficient comparison, coefficient-to-remainder
equivalence, Gevrey criterion, logarithmic optimality, action certificate,
and strong inverse theorem have self-contained proofs here. The
summability conclusions use the explicitly cited classical sectorial
characterization and generalized reconstruction theorem. Formal Lagrange
inversion, Cauchy estimates, and contraction are used in their standard
forms.

No statement is claimed Lean-verified. The repository's formal inventory
does not certify this article, and numerical testing does not upgrade its
proofs into machine-checked mathematics. The sharper exact numerical
type formulas and saddle equivalents in neighboring drafts are not needed
and have not been audited here.
The advance is specific: formal coefficient classes now control the
canonical sectorial remainders, and the subquadratic-opening regime gives
an exact directional summability classification. This is not a solution
to analytic realization in arbitrary transseries fields.

\subsection{A formalization decomposition}
The finite layer consists of the integer recurrence, partition formula,
marked positive subfamily, and log-convex weight inequalities. It needs
ordinary formal power series and finite partitions, not a new transseries
datatype. The analytic layer consists of the half-plane exponential
Taylor bound, compact normal convergence, fixed-disk contraction, and the
shrinking-disk weighted implicit theorem. A further module can package
coefficient-to-remainder equivalence. A summability layer should import
an explicit sectorial characterization rather than imply that existing
power-series APIs already prove Borel reconstruction.

\subsection{Further questions with concrete targets}
\paragraph{1. Critical negative-ray Borel summability.}
Decide ordinary negative-direction Borel summability for quadratic
feedback. The missing step is uniform control on an appropriate boundary
domain, not another coefficient-growth estimate. Inverse coordinates,
where the literal kernel is analytic for $\Rea U<0$, may help construct
a Nevanlinna--Sokal domain after returning to $q$.
% ed. (2026-09-29): the proposed method is that of the negative-ray package; the angular question is answered negatively.
\begin{ednote}
The inversion-first method proposed here is the method of \nolinkurl{Negative_Ray_Summation_Exponential_Feedback/}~\cite{ed:negativeray},
which uses it to prove fine summability in direction $\pi$ for
$\lambda_j=O((j\log j)^2)$. Ordinary (angular) summability for
$\lambda_j=j^2$ fails by \nolinkurl{Natural_Boundaries_Quadratic_Exponential_Feedback/}~\cite{ed:naturalboundaries}.
\end{ednote}

\paragraph{2. Sharp numerical remainder types.}
Compare the infimal geometric constant in strong remainders with the
coefficient type. A no-loss argument could sharpen the exponential
constant in \eqref{eq:quadraticactual} and determine whether precise type
is invariant under the inverse. Our proof establishes class equality only.
% ed. (2026-09-29): coefficient-type invariance answered by the weighted-type package.
\begin{ednote}
For coefficient types this is answered by
\nolinkurl{Sharp_Weighted_Type_Formal_Reversion/}~\cite{ed:weightedtype}:
the exact Gevrey type is invariant under the inverse (its
Theorem~\texttt{thm:inversefeedback}). The sharp constant in the remainder bounds
remains open.
\end{ednote}

\paragraph{3. Curved boundary layers.}
Since $U=q+u_2q^2+\cdots$ with $u_2>0$, the boundary $\Rea U=0$ need
not be $\Rea q=0$. Analyze the layer $\Rea q=O(|q|^2)$ and the maximal
domain of the literal canonical solution. A fixed $v$-disk does not
provide uniformity there; a curved or shrinking dependent-variable domain
may be required.

\paragraph{4. Summation and acceleration for $p\ge2$.}
The natural opening index is $p-1$, while the directly constructed germ
has opening $1$ in units of $\pi$. Investigate continuation, acceleration,
and singular barriers. Divergence of the literal positive-real kernel
does not settle analytically continued or generalized summation questions.

\paragraph{5. Signed and complex feedback.}
For $U=\sum c_jq^je^{\lambda_jU}$, find hypotheses replacing the positive
subfamily estimate. Absolute convergence can preserve the analytic disk
argument, but coefficient cancellations may conceal large kernel terms.
Identify the exact obstruction to coefficient-to-remainder equivalence.

\paragraph{6. Vector systems.}
For $U_i=\sum_jq^{a_{ij}}c_{ij}
\exp(\langle\lambda_{ij},U\rangle)$, seek a common dissipative cone,
normalize each leading component, and prove a multitype counterpart of
the bridge. The positivity argument would require a marked-tree comparison
that respects the coupled components.

\paragraph{7. Faster weights.}
The hypothesis $\mu_n^n\le D^nM_n$ excludes weights such as
$M_n=a^{n^2}$, $a>1$. Determine whether a modified shrinking-domain
argument yields a bridge in these regimes or whether the feedback model
provides counterexamples. This localizes a potential boundary of control
by formal coefficient growth.

\paragraph{8. Fully verified evaluation.}
Combine exact coefficients, interval exponential evaluation, the action
certificate, and effective weighted constants into an algorithm choosing
between action truncation and Taylor truncation. Its output should record
the method and provide a machine-checkable numerical enclosure.

\section{Conclusion}
The original coordinates really do lack an analytic dependent-variable
neighborhood when feedback slopes grow too rapidly. That failure does
not prevent uniform sectorial asymptotics. Factoring out the leading
variable exposes a fixed disk on which the exponentials are bounded;
positivity identifies its remainders with a controlled subfamily of the
formal coefficients.
This proves an exact bridge from formal regularity to canonical strong
sectorial regularity, optimal factorial--logarithmic classes, and a
complete negative-direction $k$-summability criterion for $k>1$.
The actual inverse shares the structure. Remaining problems are now
localized to sharp constants, critical aperture, angular continuation,
cancellation, and faster weights.

\appendix
\section{Reproduction and file inventory}
The source is self-contained and loads no repository notation files.
A standard TeX installation with the preamble packages suffices:
\begin{lstlisting}
pdflatex -interaction=nonstopmode -halt-on-error article.tex
pdflatex -interaction=nonstopmode -halt-on-error article.tex
pdflatex -interaction=nonstopmode -halt-on-error article.tex
\end{lstlisting}
The exact checks require only Python's standard library; optional numerical
checks use mpmath:
\begin{lstlisting}
python code/verify.py --order 120 --exact-only
python -m pip install mpmath
python code/verify.py --order 120
\end{lstlisting}
The full run regenerates \path{data/verification.json}, four exact
coefficient CSVs, and the sectorial diagnostics CSV. The CSV column
\texttt{n\_factorial\_times\_u\_n} contains $n!u_n$, not $u_n$.
Numerator and denominator columns give the ordinary coefficients exactly.
The package also contains a README, source/provenance notes, proof-status
summary, recorded output, and a Makefile. No repository changes were made
and no external font files are distributed.

\clearpage
\begingroup\raggedright
\begin{thebibliography}{99}
\bibitem{repo}
V. Reshetnikov, \emph{ProveIt}, transseries documentation README,
snapshot \texttt{0c973d8f5e7ef4550f4df1bf486d890037fd9f14},
consulted 29 September 2026.
\url{https://github.com/VladimirReshetnikov/ProveIt/tree/0c973d8f5e7ef4550f4df1bf486d890037fd9f14/Analysis/Transseries/docs/series-and-transseries}.

\bibitem{feedback}
\emph{A sharp regularity classification for countable exponential-feedback
transseries}, research draft prepared for V. Reshetnikov, 29 September
2026. In \cite{repo}, directory
\path{Exponential_Feedback_Regularity_Classification/},
\path{article.tex} and \path{README.md}. Its left-sector realization
proof explicitly does not claim uniform Gevrey remainders.

\bibitem{finitecore}
\emph{Finite-Core Universality and Sharp Large Order: Countable
exponential-feedback transseries}, research draft prepared for
V. Reshetnikov, 29 September 2026. In \cite{repo}, directory
\path{Finite_Core_Universality_Exponential_Feedback/}; README consulted
for scope and non-duplication.

\bibitem{nikolaev}
N. Nikolaev, \emph{Gevrey asymptotic implicit function theorem},
L'Enseignement Math\'ematique \textbf{70} (2024), no.~1/2, 251--282.
\url{https://doi.org/10.4171/LEM/1061}.
Preprint: \url{https://arxiv.org/abs/2112.08792}.

\bibitem{lms}
A. Lastra, S. Malek, and J. Sanz,
\emph{Summability in general Carleman ultraholomorphic classes},
Journal of Mathematical Analysis and Applications \textbf{430} (2015),
no.~2, 1175--1206.
\url{https://doi.org/10.1016/j.jmaa.2015.05.046}.
Preprint: \url{https://arxiv.org/abs/1402.1669}.
We use the preprint's Example 2.2, Definition 3.3, Theorem 3.18,
and Remark 3.19(i), with ordinary coefficient notation.

\bibitem{sanz}
J. Sanz, \emph{Flat functions in Carleman ultraholomorphic classes via
proximate orders}, Journal of Mathematical Analysis and Applications
\textbf{415} (2014), no.~2, 623--643.
\url{https://doi.org/10.1016/j.jmaa.2014.01.083}.
Preprint: \url{https://arxiv.org/abs/1402.2627}.

\bibitem{sokal}
A. D. Sokal, \emph{An improvement of Watson's theorem on Borel
summability}, Journal of Mathematical Physics \textbf{21} (1980),
no.~2, 261--263. \url{https://doi.org/10.1063/1.524408}.
Cited as a possible boundary-domain route, not an already verified
hypothesis for this model.
% ed. (2026-09-29): three bibliography entries added for the editorial notes.
\bibitem{ed:negativeray}
[Editorial addition, ProveIt, 2026-09-29.] \emph{Negative-Ray Summation of
Countable Exponential-Feedback Transseries}, research article prepared for
V. Reshetnikov, 29 September 2026. In \cite{repo}, directory
\path{Negative_Ray_Summation_Exponential_Feedback/}; present at the
snapshot used by this article.
\bibitem{ed:naturalboundaries}
[Editorial addition, ProveIt, 2026-09-29.] \emph{Natural Boundaries of
Quadratic Exponential Feedback}, research draft prepared for
V. Reshetnikov, 29 September 2026. Directory
\path{Natural_Boundaries_Quadratic_Exponential_Feedback/} of the same
series; filed after this article.
\bibitem{ed:weightedtype}
[Editorial addition, ProveIt, 2026-09-29.] \emph{Sharp Weighted Type Under
Formal Reversion}, research article prepared for V. Reshetnikov,
29 September 2026. Directory \path{Sharp_Weighted_Type_Formal_Reversion/}
of the same series; filed with this article.
\end{thebibliography}
\endgroup
\end{document}
